\subsubsection{Retour persistant des cinq régions}

Soit
\[
\mathcal F_\pi=\Psi^{-1}(010211)\subset\mathcal U_6
\]
la microfibre relative au catalogue. Rappelons ses cinq éléments, maintenant séparés en bloc initial et bloc terminal :
\[
\begin{aligned}
 &(501,614,498\mid169,272,272),\\
 &(810,923,870\mid169,272,272),\\
 &(810,983,810\mid169,272,272),\\
 &(870,923,810\mid169,272,272),\\
 &(870,923,810\mid418,674,521).
\end{aligned}
\]
Nous définissons la projection de retour
\[
p_{\rm ret}:\mathbb Z^6\longrightarrow\mathbb Z^3,
\qquad
p_{\rm ret}(u_1,\ldots,u_6)=(u_4,u_5,u_6).
\]

\begin{proposition}[Bloc terminal persistant]
\label{p12-83-c13-s3:prop:bloque-terminal-persistente}
Le multiensemble \(p_{\rm ret}(\mathcal F_\pi)\) possède un unique élément de
multiplicité maximale :
\[
b_\pi=(169,272,272),
\qquad
\operatorname{mult}_{\mathcal F_\pi}(b_\pi)=4.
\]
L'unique bloc terminal restant est
\[
b_\pi^-=(418,674,521)
\]
et a une multiplicité de un.
\end{proposition}

\begin{proof}
L'affirmation s'obtient par énumération exacte des cinq lignes du catalogue qui satisfont \(\Psi(U)=010211\). Quatre lignes possèdent le bloc terminal \((169,272,272)\) ; la cinquième, qui porte le retour muté, possède \((418,674,521)\). La multiplicité maximale est quatre et n'est atteinte qu'une seule fois.
\end{proof}

La différence orientée entre le retour muté et le retour persistant est
\[
 \omega_\pi=b_\pi^- - b_\pi=(249,402,249),
 \qquad
 \langle(1,1,1),\omega_\pi\rangle=900.
\]
Cette identité sépare deux objets qu'il ne faut pas confondre. L'entier
\(169\) est une coordonnée d'une donnée terminale, c'est-à-dire une donnée de degré
zéro. Le vecteur \(\omega_\pi\) peut être lu comme une cochaîne de degré un
après orientation de l'arête qui relie les deux blocs. L'appeler cocycle exigerait,
en outre, de fixer un complexe de cochaînes et de vérifier que son cobord s'annule.
La prolongation analytique construite plus loin ne présuppose pas cette
identification.

La sélection de \(169\) n'exige pas de privilégier la première position du bloc.
Le groupe symétrique \(\mathfrak S_3\) agit en permutant ses trois coordonnées.
Le stabilisateur de \(b_\pi\) échange les deux occurrences de \(272\) et
fixe l'unique indice de coordonnée dont la valeur a une multiplicité égale à un.
Pour un bloc de type \((r,s,s)\), avec \(r\ne s\), notons cet indice
\(\sigma(b)\) et sa valeur sélectionnée
\[
 \operatorname{sing}(b)=b_{\sigma(b)}.
\]
Pour \(\tau\in\mathfrak S_3\), le sélecteur d'indices est équivariant et la valeur sélectionnée est invariante :
\[
 \sigma(\tau\cdot b)=\tau\bigl(\sigma(b)\bigr),\qquad
 \operatorname{sing}(\tau\cdot b)=\operatorname{sing}(b).
\]
De manière équivalente,
\[
\operatorname{sing}(b)
=\sum_{j=1}^3b_j-2\operatorname{moda}(b).
\]
En particulier,
\[
\boxed{\operatorname{sing}(b_\pi)=169.}
\]

\begin{definition}[Sélecteur résiduel de la microfibre]
\label{p12-83-c13-s3:def:selector-residual-pi}
Le sélecteur \(\rho_\pi\) applique successivement deux opérations :
\begin{enumerate}
  \item sélectionne le bloc terminal de multiplicité maximale au sein de \(p_{\rm ret}(\mathcal F_\pi)\) ;
  \item sélectionne l'indice de coordonnée fixé par le stabilisateur de
  ce bloc et retient la valeur correspondante.
\end{enumerate}
D'après la proposition~\ref{p12-83-c13-s3:prop:bloque-terminal-persistente},
\[
\boxed{\rho_\pi=169.}
\]
\end{definition}

\begin{theorem}[Unicité équivariante du retour]
\label{p12-83-c13-s3:thm:unicidad-retorno-169}
Parmi les sélecteurs qui
\begin{enumerate}
  \item sont invariants sous les permutations des cinq régions ;
  \item sélectionnent un indice de coordonnée de manière équivariante sous les permutations des trois coordonnées terminales ;
  \item sélectionnent le bloc de multiplicité maximale ;
  \item sélectionnent un indice de coordonnée fixe par son stabilisateur,
\end{enumerate}
la valeur de retour est déterminée de manière unique et vaut \(169\).
\end{theorem}

\begin{proof}
L'invariance sous \(\mathfrak S_5\) empêche d'utiliser l'ordre des lignes. La
troisième condition sélectionne \(b_\pi\), qui est l'unique mode d'après la
proposition~\ref{p12-83-c13-s3:prop:bloque-terminal-persistente}. Son stabilisateur dans
\(\mathfrak S_3\) possède deux orbites sur les indices de coordonnées :
la paire qui porte la valeur \(272\) et l'indice ponctuel qui porte
\(169\). La quatrième condition sélectionne la seconde orbite. L'
équivariance transporte l'indice lors de la permutation des coordonnées ; la
valeur portée par cet indice reste égale à \(169\).
\end{proof}

Cet argument améliore la lecture positionnelle du retour. Le \(169\) n'est pas « la première coordonnée » d'une ligne choisie : c'est un invariant du multiensemble régional et de l'action de ses symétries de réordonnancement. Il n'est pas non plus sélectionné pour sa rareté globale. Dans les \(468\) émissions du catalogue, \(169\) apparaît \(84\) fois, uniformément réparti avec \(14\) occurrences dans chacune des six positions. L'unicité du théorème~\ref{p12-83-c13-s3:thm:unicidad-retorno-169} est donc relationnelle et régionale : elle réside dans la combinaison de microfibre, de persistance modale et de stabilisateur, et non dans la fréquence isolée de l'entier.

\subsubsection{Longueur cylindrique et degré analytique}

Soit \(\mathcal P\) le module libre engendré par les mots régionaux et
soit \(F^n\mathcal P\) sa filtration par longueur. La graduation associée
enregistre le premier niveau auquel apparaît chaque mot. Pour un paramètre
formel \(z\), le caractère ordinaire de longueur est
\[
\chi_z([w])=z^{|w|}.
\]
La multiplicativité
\[
\chi_z([uv])=\chi_z([u])\chi_z([v])
\]
exprime que la concaténation additionne les longueurs. Un mot de longueur
\(n\) a pour degré formel \(n\), représenté par le monôme \(z^n\).
L'évaluation ultérieure en \(z=\alpha\) attribue à ce monôme le poids
numérique \(\alpha^n\).

\begin{definition}[Compatibilité de degré]
\label{p12-83-c13-s3:def:compatibilidad-grado}
Le lecteur analytique régional est compatible avec la filtration cylindrique
lorsqu'il envoie la composante homogène de longueur \(n\) sur la composante
formelle homogène de degré \(n\) :
\[
\operatorname{gr}_n\mathcal P\longrightarrow
\mathbb R z^n.
\]
L'évaluation en \(z=\alpha\) est une opération ultérieure sur
l'expression formelle.
\end{definition}

Pour tout mot \(w\in\mathcal F_\pi\subset\mathcal U_6\), la donnée commune de retour apparaît au degré six :
\[
\rho_\pi\chi_z([w])=169z^6
\xrightarrow{\;z=\alpha\;}169\alpha^6.
\]
Ce terme enregistre une seule fois l'impulsion commune de la microfibre.
Ce six est la longueur du mot régional. Ce n'est pas la longueur des
marches fermées de la section~\ref{regcl27:seccion}. Ces marches parcourent
simultanément les quatre points d'ancrage
\(U_\pi,U_e,U_\varphi,U_{\rm APP}\) et ont une longueur minimale de huit. Un
mot, une arête du graphe de blocs et une marche fermée sont des objets de
types différents ; la graduation utilisée ici appartient au premier.

\subsubsection{Le caractère déterminantal du transport commun}

La fermeture dodécaphasique fournit \(\alpha\). La complétion des trois
couples nonadiques fournit l'échelle \(10^3\), de sorte que
\[
A=10^3\alpha.
\]
La coordinatisation angulaire archimédienne étant choisie, la coordonnée commune en radians est
\[
x=\frac{\pi A}{180}=\frac{50\pi}{9}\alpha.
\]
Soit \(R\) une involution réelle de trace nulle en dimension deux,
\[
R^2=I_2,
\qquad
\operatorname{tr}R=0,
\]
et considérons le transport formel bicouche
\[
\mathcal T(x,y)=\exp(-xI_2-yR).
\]
La variable orientée \(y\) n'a pas encore été évaluée. Cependant, l'action
de \(\mathcal T\) sur la droite déterminante en est indépendante :
\[
\begin{aligned}
\det\mathcal T(x,y)
&=\exp\!\left(\operatorname{tr}(-xI_2-yR)\right)\\
&=e^{-2x}.
\end{aligned}
\]
On définit
\[
\boxed{
D_A=e^{-2x}=e^{-100\pi\alpha/9}.}
\]
L'indice rappelle qu'il s'agit du déterminant de la contraction commune associée à \(A\), et non d'une fonction de \(C_*^{\rm HMT}\).

\begin{proposition}[Caractère de la droite déterminante]
\label{p12-83-c13-s3:prop:caracter-determinante}
L'application
\[
\delta_A:(\mathbb N_0,+)\longrightarrow(\mathbb R_{>0},\cdot),
\qquad
\delta_A(k)=D_A^k,
\]
est l'unique caractère multiplicatif dont la valeur sur une prolongation
élémentaire est \(D_A\).
\end{proposition}

\begin{proof}
La multiplicativité donne
\[
\delta_A(k)
=\delta_A(\underbrace{1+\cdots+1}_{k\text{ fois}})
=\delta_A(1)^k=D_A^k.
\]
La même identité prouve l'existence et l'unicité.
\end{proof}

Le déterminant reste invariant par conjugaison de \(\mathcal T\) et par échange des feuillets \(R\mapsto-R\). La queue construite à partir de \(D_A\) appartient donc au mode commun. L'orientation apparaîtra dans le signe avec lequel le caractère régional corrige la phase.

\Needspace{29\baselineskip}
\subsubsection{Règles du lecteur analytique régional}

L'extension analytique est fixée par les règles suivantes.

\begin{definition}[Lecteur régional déterminantal]
\label{p12-83-c13-s3:def:lector-regional-determinantal}
Le lecteur régional déterminantal satisfait :
\begin{enumerate}
  \item \emph{compatibilité de degré} : la longueur cylindrique est le degré
  analytique formel ;
  \item \emph{décomposition de renouvellement} : l'impulsion finie de retour et
  la continuation stricte appartiennent à des termes additifs distincts ;
  \item \emph{normalisation de continuation} : après l'enregistrement, une seule
  fois, du résidu \(\rho_\pi\), la branche survivante commence à l'unité de
  la droite déterminant ;
  \item \emph{covariance par concaténation} : chaque prolongation élémentaire
  agit au moyen de
  \(\bigwedge^2\mathcal T=D_A\,\operatorname{Id}_{\bigwedge^2\mathbb R^2}\) ;
  \item \emph{orientation} : dans la convention cardinale fixée pour APP, le retour régional appartient au secteur compensateur et est soustrait de la phase d'action de cinquième degré.
\end{enumerate}
\end{definition}

La troisième règle a un contenu mathématique. La valeur \(169\) est enregistrée
une fois comme amplitude du retour ; elle n'est pas multipliée à nouveau à chaque
prolongation. La continuation stricte compte l'unique droite qui se poursuit et
la normalise par son unité. C'est la différence entre une récurrence de
renouvellement et l'itération homogène de l'impulsion initiale.

La nécessité de déclarer cette normalisation peut être démontrée. Pour tout
\(\lambda\in\mathbb R\), la collection
\[
F_\lambda(z)
=169z^6+\lambda\frac{D_Az^7}{1-D_Az}
\]
conserve le même retour fini et la même récurrence multiplicative des
coefficients de la queue. Le catalogue et cette récurrence ne distinguent pas à eux seuls
\(\lambda\). L'unité de la droite déterminant fixe
\(\lambda=1\) avant l'évaluation de \(z=\alpha\). Ainsi, la normalisation
ne s'obtient pas en ajustant la valeur finale ; elle fait partie de la définition du
caractère.
